\section{Pricing and valuing forward commitments}

These entries price forwards, futures and swaps by arbitrage and replication and value them during their lives: equity, interest-rate and bond forwards and futures, and interest-rate, currency and equity swaps. The cost-of-carry formulas with income, yield and storage are in \hyperref[res:cost-of-carry-forward-price]{the cost of carry}; FRA and swap valuation with separate projection and discounting curves is in \hyperref[mod:multi-curve-fra-swap-valuation]{multi-curve valuation}. Unless an example says otherwise, it uses one curve for both projection and discounting, a simplification that matches pre-2008 textbook practice and today's OIS and SOFR swaps, not term-rate swaps. Numbers are checked in \texttt{notes/artefacts/forward-swap-valuation/forward-swap-checks.py}; rates are continuously compounded unless stated.

\begin{qfmethod}[Arbitrage, replication and the cost of carry: why a forward commitment is priced by the cost of a replicating position rather than by expected prices, and the costs and benefits that move the forward price]\label{meth:arbitrage-replication-carry}
\qffield{Statement} \emph{Arbitrage} locks in a riskless profit from simultaneous trades, and pricing assumes that ``no arbitrage opportunities exist'' \citep[chap.~1, sec.~1.9, p.~39]{s7-hulloptions}. Two portfolios with the same future payoffs must therefore cost the same today (the law of one price). A long forward on an investment asset is \emph{replicated} by buying the asset with borrowed money: at maturity the buyer owns the asset and owes the loan with interest, the same position as a long forward struck at that amount. Hence $F_0=S_0e^{rT}$ without income, $(S_0-I)e^{rT}$ with income of present value $I$, $S_0e^{(r-q)T}$ with yield $q$ (chap.~5, secs.~5.4--5.6, pp.~127--132). The \emph{cost of carry}, financing plus storage less income and, for consumption assets, convenience yield, is what moves $F_0$ away from $S_0$ (sec.~5.12, pp.~143--144; \hyperref[res:cost-of-carry-forward-price]{the cost of carry}).

Because the replication uses no forecast, the forward price does not depend on the expected future spot price or on investors' risk preferences: the same price would be found in a world in which investors are risk-neutral, discounting expected payoffs at the risk-free rate, which is why that device can be used for pricing (risk-neutral valuation, Hull, sec.~13.2, pp.~292--293).

\qffield{Practice} \emph{Original} (spot 100, $r=4$ percent, one year, no income): $F_0=104.0811$. If a dealer quotes 106, borrow 100, buy the asset and sell forward: at maturity deliver for 106, repay 104.0811, profit 1.9189. At 102, short the asset, lend the proceeds and buy forward: profit 2.0811. If investors expect 8 percent growth, $\mathbb{E}[S_T]=108.33$, but the forward is still 104.0811: the gap is the risk premium, not an arbitrage. \emph{Reproduced} (Hull): $40e^{0.05\times0.25}=40.50$ (Example~5.1, p.~128); $(900-40e^{-0.03/3})e^{0.04\times0.75}=886.60$ (Example~5.2, p.~131).

\emph{Commitments against contingent claims.} A forward commitment's linear payoff is replicated by a static position, bought once and held, so arbitrage fixes one price. An option's non-linear payoff is not: static arbitrage gives only bounds and, for European options, put--call parity, $c+Ke^{-rT}=p+S_0$ (Hull, secs.~11.3--11.4, pp.~252--257). A single option price needs a model of how the underlying can move; in Hull's one-step binomial tree a position of $\Delta$ shares against a short option is riskless and must earn the risk-free rate, which fixes the option price, and over many steps the position has to be rebalanced (sec.~13.1, pp.~288--291).

\qffield{Conditions} Hull's assumptions: some market participants trade without costs, pay one tax rate, borrow and lend at the risk-free rate and exploit arbitrage (sec.~5.3, p.~126). With transaction costs, borrow fees or different borrowing and lending rates the price becomes a band. For a consumption asset only the upper bound is enforced (\hyperref[res:cost-of-carry-forward-price]{consumption assets}).

\qffield{Pitfalls} Reading the forward price as a forecast. Treating an arbitrage that needs a short sale as free when borrowing is costly or recalled (\hyperref[def:securities-lending-fee-rebate]{lending fees}).

\qffield{See also} \hyperref[meth:forward-price-value-lifecycle]{price and value over the life}; \hyperref[def:derivative-uses-risks]{arbitrageurs}.

\end{qfmethod}

\begin{qfmethod}[Price and value of a forward over its life: zero value at inception, value to long and short before maturity, and the payoff at maturity]\label{meth:forward-price-value-lifecycle}
\qffield{Statement} The \emph{price} of a forward is the delivery price $K$ fixed in the contract; the \emph{value} is what the position is worth. At inception $K$ is set to the forward price, so the contract is worth zero to both sides. Later, with the current forward price $F_t$ for the same delivery date, the long is worth $f_t=(F_t-K)e^{-r(T-t)}$, the present value of the gain from selling forward at $F_t$ what it has agreed to buy at $K$; the short holds $-f_t$ \citep[chap.~5, sec.~5.7, pp.~133--134]{s7-hulloptions}. For an asset without income, $f_t=S_t-Ke^{-r(T-t)}$; at maturity the value is the payoff $S_T-K$.

\qffield{Practice} \emph{Reproduced}: $f=(26.28-24)e^{-0.05}=2.17$ (Hull, Example~5.4, p.~133). \emph{Original}, continuing the 104.0811 forward: six months later with the share at 108, the long is worth $108-104.0811e^{-0.02}=5.9799$ and the short $-5.9799$; at maturity with $S_T=110$ the long receives 5.9189. The value moves with spot, rates and, below, income.

\qffield{Conditions} The formula is the same for investment and consumption assets, because it uses only the forward price observed at $t$ (Hull, p.~133). It values a forward, not a futures contract, whose daily settlement resets its value to zero each day.

\qffield{Pitfalls} Confusing the forward price, which is quoted, with the forward's value, which is zero only at inception. Discounting the price difference to the valuation date rather than from the delivery date.

\qffield{See also} \hyperref[meth:equity-forward-futures-valuation]{equity forwards}; \hyperref[meth:rate-bond-forward-valuation]{rate and bond forwards}.

\end{qfmethod}

\begin{qfresult}[Forward and futures prices: equal when interest rates are known, different when the asset is correlated with rates, and the practical reasons they diverge]\label{res:forward-futures-price-difference}
\qffield{Statement} When the short-term risk-free rate is constant, or a known function of time, the forward and futures prices for the same delivery date are equal in theory; with unpredictable rates they are not. If the asset is positively correlated with rates, a long futures position gains when rates have risen, reinvesting the daily settlement at higher rates, and loses when rates have fallen, financing the loss at lower rates, so futures prices tend to exceed forward prices; with negative correlation forward prices tend to be higher \citep[chap.~5, sec.~5.8, p.~135]{s7-hulloptions}. Hull judges the differences for contracts of a few months ``sufficiently small to be ignored'' in most circumstances, and lists taxes, transaction costs, margin, the clearing house's reduction of default risk and differences in liquidity as practical reasons for divergence.

\qffield{Practice} \emph{Original checks.} With a deterministic but time-varying short rate, a futures position \emph{tailed} to $e^{-\int_{t+1}^{T}r}$ contracts and with each day's settlement invested to $T$ ends with exactly $S_T-F_0$, the forward payoff, on a simulated path: forward and futures prices coincide. In a two-period tree with short rate 4 percent, then 6 or 2 percent with equal risk-neutral probability (simple compounding), and the asset at 110 when rates rise and 92 when they fall, the futures price is 105.2200 and the forward price 105.0012; with the asset's moves reversed the futures price falls to 104.8600 while the forward stays at 105.0012. Eurodollar and SOFR futures, whose payoff is linked to rates themselves, need a \emph{convexity adjustment} for this reason (Hull, sec.~6.3, pp.~163--164).

\qffield{Conditions} The argument needs the same delivery and settlement terms; delivery options and different underlying grades also separate futures from forward prices (\hyperref[meth:rate-bond-forward-valuation]{bond futures}).

\qffield{Pitfalls} Applying the correlation argument to the forward's value rather than the price; using futures prices as forward prices for long-dated rate contracts without the convexity adjustment.

\qffield{See also} \hyperref[mech:futures-margin-settlement]{daily settlement}; \hyperref[def:futures-hedge-basis]{tailing a hedge}.

\end{qfresult}

\begin{qfdefinition}[Swaps compared with forwards: a swap as a strip of forward contracts and as a package of bond positions, off-market forwards, and how the value of each exchange changes]\label{def:swaps-as-forward-strips}
\qffield{Statement} A swap ``can be viewed as a package of'' forward rate contracts, or ``as a package of cash flows from buying and selling cash market instruments'': a fixed-rate payer is long a floating-rate bond and short a fixed-rate bond \citep[chap.~64, pp.~1577--1578]{s7-fabozzifixed}. Swaps differ from a strip of forwards in that they run longer than liquid forwards and futures and are executed in one transaction. Every exchange in a swap is at the same fixed rate, so each is an \emph{off-market} forward with non-zero value; the fixed rate is chosen so that their values sum to zero \citep[chap.~7, sec.~7.7, pp.~185--186]{s7-hulloptions}. When the term structure slopes upward and the fixed rate is received, early exchanges have positive and later ones negative value, so the swap is expected to become negative over time (Figure~7.9, p.~186). Futures differ from both in being standardised and settled daily. \emph{Price and value.} A swap's price is its fixed rate, set at inception so that the swap's value is zero; afterwards the rate stays fixed while the value moves with the curve, exactly as a forward's delivery price stays fixed while its value moves (\hyperref[meth:forward-price-value-lifecycle]{forward price and value}).

\qffield{Practice} \emph{Original} (annual OIS par swap rates 3.60, 3.50, 3.55, 3.65 and 3.75 percent; one curve for projection and discounting): bootstrapped one-year forwards are 3.60, 3.40, 3.66, 3.97 and 4.19 percent, and the five-year par swap rate $(1-P_5)/\sum_iP_i$ recovers 3.75 percent. Each exchange against 3.75 percent has non-zero value; to the floating receiver the first two are negative and the last positive, and the five sum to zero.

\qffield{Conditions} The forward-strip view needs forwards on the swap's own floating index, settled at its payment dates; with separate projection and discounting curves, each exchange is valued as in \hyperref[mod:multi-curve-fra-swap-valuation]{multi-curve valuation}. The bond-package view needs a floater worth par on reset dates, true on one curve.

\qffield{Pitfalls} Treating a par swap as a series of zero-value forwards. Assuming the swap's value stays near zero because it started there.

\qffield{See also} \hyperref[meth:swap-valuation-replication]{swap valuation}; \hyperref[def:interest-rate-swaps]{interest-rate swaps}.

\end{qfdefinition}

\begin{qfmethod}[Equity forwards and futures: pricing with discrete dividends and with a dividend yield, replication, index futures, and the value of an existing equity forward]\label{meth:equity-forward-futures-valuation}
\qffield{Statement} For a share paying known dividends with present value $I$ before delivery, $F_0=(S_0-I)e^{rT}$. Replication borrows the full $S_0$, buys the share and pays each dividend to the lender when it is received; the debt left at $T$ is $S_0e^{rT}$ less each dividend grown to $T$, which is $(S_0-I)e^{rT}$, so the position is a long forward at $F_0$. Borrowing only $S_0-I$ would leave the purchase short of funds by $I$ at the start. For an index with continuous dividend yield $q$, $F_0=S_0e^{(r-q)T}$, replicated by holding $e^{-qT}$ units and reinvesting dividends \citep[chap.~5, secs.~5.5--5.6 and~5.9, pp.~130--132 and~135--137]{s7-hulloptions}. An existing long forward is worth $S_t-I_t-Ke^{-r(T-t)}$, with $I_t$ the present value of the dividends still to come. If the futures price exceeds $S_0e^{(r-q)T}$, index arbitrage buys the index shares and sells futures; otherwise it does the reverse (pp.~136--137). Equity futures prices equal forward prices to Hull's approximation (\hyperref[res:forward-futures-price-difference]{forwards and futures}).

\qffield{Practice} \emph{Original} (spot 100, 4 percent, one year): dividends of 1.50 in three and nine months have present value 2.940743 and $F_0=101.020320$. The replication is self-financing date by date: at the start 100 is borrowed and spent on the share; at three and nine months each 1.50 dividend repays part of the loan; at maturity the share is held against a debt of 101.020320. If a dealer quoted 102, running this position and selling forward would earn 0.979680 at maturity and nothing before. Six months later, at 108 with one dividend left, a long struck at 101.0203 is worth 7.4949. With a 2 percent continuous yield instead, $F_0=102.0201$. An index at 4{,}000 with a 2 percent yield and 4 percent rates has a three-month futures price of 4{,}020.05. \emph{Reproduced}: $(50-2.162)e^{0.08\times10/12}=51.14$ (Hull, Example~5.3, p.~132).

\qffield{Conditions} Known dividends; with uncertain dividends the forward price embeds the market's dividend forecast, which dividend futures and swaps trade separately. Index arbitrage needs the basket to be traded at the index level; it failed in October~1987 (\hyperref[res:cost-of-carry-forward-price]{Business Snapshot~5.4}).

\qffield{Pitfalls} Subtracting dividends paid after delivery; treating a yield applied to the price as the same as fixed cash dividends; leaving out short-sale and borrow costs in the reverse arbitrage.

\qffield{See also} \hyperref[meth:swap-valuation-replication]{equity swaps}; \hyperref[meth:forward-price-value-lifecycle]{forward value}.

\end{qfmethod}

\begin{qfmethod}[Interest-rate and bond forwards and futures: FRA rate and value, SOFR futures hedges, bond forward prices with coupons and accrued interest, conversion factors and the cheapest-to-deliver bond]\label{meth:rate-bond-forward-valuation}
\qffield{Statement} \emph{FRA.} A forward rate agreement fixes a rate $K$ for a future period on a notional. On one curve the FRA rate is the forward rate $F=(P_1/P_2-1)/\tau$, and a contract paying $K$ and receiving the floating rate is worth $N\tau(F-K)P_2$ if settled at the period end \citep[chap.~4, sec.~4.9, pp.~110--112]{s7-hulloptions}; start settlement and separate curves are in \hyperref[mod:multi-curve-fra-swap-valuation]{multi-curve valuation}. \emph{Uses} (Interpretation): a company that will borrow at a floating rate later pays fixed in an FRA to lock in its borrowing rate, and a lender or investor in floating-rate assets receives fixed to lock in its return; strips of FRAs or rate futures do the same for several periods, which a swap packages (\hyperref[def:swaps-as-forward-strips]{swaps as forward strips}). \emph{SOFR futures} are quoted as $100-R$, where $R$ compounds one-day SOFR over the contract's three months, and pay \$25 per basis point; a borrower locks in a rate by selling them (sec.~6.3, pp.~161--163).

\emph{Bond forwards.} With full (cash) price $B_0$ and coupons of present value $I$ before delivery, the forward full price is $(B_0-I)e^{rT}$, and the quoted (clean) forward price subtracts the accrued interest at delivery (cash price equals quoted price plus accrued interest, sec.~6.1, p.~154; sec.~6.2, pp.~158--160). \emph{Treasury bond futures} let the short choose which eligible bond to deliver; the long pays the futures settlement price times the bond's \emph{conversion factor} plus accrued interest. The conversion factor is the price per dollar the bond would have if all rates were 6 percent with semiannual compounding, maturity rounded down to the nearest three months (pp.~156--157; CBOT Rule 18101.B states the 6 percent basis \citep[Rule 18101.B]{cbot-treasury-chapters-2026}). The short delivers the \emph{cheapest-to-deliver} bond, minimising quoted price less settlement price times conversion factor (p.~158).

\qffield{Practice} \emph{Reproduced} (Hull): conversion factors 1.4623 for a 10 percent bond with 20 years and 1.2199 for an 8 percent bond with 18 years and 3 months (p.~157); delivery costs 2.69, 1.87 and 2.12 make the second bond cheapest (Example~6.1, p.~158); for the 12 percent bond of Example~6.2 the cash price 116.978, coupon present value 5.803, cash forward 119.711 and quoted forward 114.859 give a futures price of 71.79 with a conversion factor of 1.6 (pp.~159--160).

\emph{Example~6.3} (pp.~162--163). A borrower will pay three-month SOFR plus 200 basis points on \$100 million for three months from 16 December 2021, and Hull accrues both the loan and the futures on a quarter year ($\tau=0.25$). One contract moves $\$1\text{ million}\times0.0001\times0.25=\$25$ per basis point, so 100 contracts match the loan's $N\tau$. The quote 99.990 implies a reference rate of $100-99.990=0.01$ percent. A short position gains $N\tau(R-0.01\%)$ when SOFR settles at $R$, and the loan costs $N\tau(R+2\%)$, so the net cost is $N\tau(0.01\%+2\%)$, 2.01 percent, whatever $R$ turns out to be.

If the contract settles at 99.200 ($R=0.80$ percent), the loan pays 2.80 percent, $0.25\times\$100\text{ million}\times0.028=\$700{,}000$, and the futures gain 79 basis points, $100\times\$25\times79=\$197{,}500$: net \$502{,}500. If it settles at 100.400 ($R=-0.40$ percent), the loan pays 1.60 percent, \$400{,}000, and the futures lose 41 basis points, \$102{,}500: net \$502{,}500 again. The lock is exact only in this simplification: the futures' reference quarter runs from 15 December 2021 to 16 March 2022 (91 days) while the loan's three months from 16 December have 90, gains are settled daily rather than at the end, and a 92-day act/360 loan would need about 102 contracts.

\emph{Original.} From discount factors 0.990 and 0.979 at three and six months the 3$\times$6 FRA rate is 4.4450 percent (act/360, 91 days); an existing FRA paying 4.20 percent on 10 million is worth $+6{,}063$ to the fixed payer. A 5 percent bond with clean price 101.20 and accrued interest 0.90 (full 102.10), a 2.50 coupon in 0.35 years and repo at 4.2 percent has a six-month forward full price of 101.7510 and a quoted forward of 101.3760 after 0.375 of accrued interest; if the forward full price later rises 0.40 with 0.25 years left, the long is worth $0.3958$.

\qffield{Conditions} Hull's futures price from the cheapest-to-deliver bond ignores the short's delivery options, whose value makes the futures price ``lower than it would be without the option'', and the effect of daily settlement (pp.~158--159). Conversion factor rules are CBOT's as retrieved in October~2026. SOFR futures replaced Eurodollar futures after LIBOR ended.

\qffield{Pitfalls} Hull's text of Example~6.3 subtracts \$197{,}000 rather than \$197{,}500, and in the negative-rate case computes the 41 basis point loss as \$102{,}250, not \$102{,}500, and totals \$502{,}250; both totals are \$502{,}500. Dropping the borrowing spread: the futures lock only the SOFR component, 0.01 percent here, and the spread is paid on top. Using a clean price where the full price belongs in the carry formula; ignoring which bond is cheapest when hedging with bond futures.

\qffield{See also} \hyperref[mod:forward-rates-fra]{forward rates and FRAs}; \hyperref[meth:spot-par-forward-curves]{forward curves}.

\end{qfmethod}

\begin{qfmethod}[Valuing swaps by replication: interest-rate swaps by bonds and by forward rates, fixed-for-fixed currency swaps, and equity swaps, at inception and during their life]\label{meth:swap-valuation-replication}
\qffield{Statement} \emph{Interest-rate swap.} Receive-fixed is long a fixed-rate bond and short a floating-rate bond: $V=B_{\text{fix}}-B_{\text{fl}}$, with the floater worth par on a reset date on one curve. Equivalently it is the sum of the exchanges valued as FRAs at today's forwards, \citep[chap.~7, sec.~7.6, pp.~183--185]{s7-hulloptions}; for the bond package see \citep[chap.~64, pp.~1577--1578]{s7-fabozzifixed}. Interpretation: on one curve both equal the remaining annuity times the difference between the contract's fixed rate and today's par swap rate for the remaining term. At inception the par rate is $(1-P_n)/\sum_iP_i$ and the swap is worth zero.

\emph{Fixed-for-fixed currency swap.} Principals are exchanged at the start and end; the swap is a long bond in one currency and a short bond in the other, $V=SB_F-B_D$ in domestic currency with spot $S$ in domestic per foreign, or a portfolio of currency forwards (Hull, secs.~7.8--7.9, pp.~186--192). \emph{Equity swap.} The equity leg pays the total return on a notional each period; borrowing the notional at each payment date and investing it in the index, dividends reinvested, replicates it, so an equity-for-floating swap is worth zero at inception and just after each payment date (chap.~34, sec.~34.4, pp.~777--778). Take a total-return index, one curve for projection and discounting, and payment at each reset. The next payment, $N(E_{\text{next}}/E_0-1)$ with $E_0$ the index at the last reset, is replicated by holding $NE_t/E_0$ in the index and owing $N$ at the payment date, so it is worth $N\bigl(E_t/E_0-P(t,T_{\text{next}})\bigr)$. Each later payment is worth $N(P_{k-1}-P_k)$; they telescope, and the whole equity leg is worth $N(E_t/E_0-P_n)$, with $P_n$ the discount factor to the last payment.

The floating leg mirrors it. With compounded overnight SOFR, let $A_t$ be the growth already realised since the last reset; investing $NA_t$ overnight until the payment date replicates the next floating payment, worth $N(A_t-P(t,T_{\text{next}}))$, and later payments are worth $N(P_{k-1}-P_k)$ as on the equity leg. The net equity-for-SOFR value to the equity receiver is therefore $N(E_t/E_0-A_t)$, less the present value of any spread on the floating side. With a rate fixed at the last reset, $R_0$ for $\tau_0$, the next floating payment is worth $NR_0\tau_0P(t,T_{\text{next}})$ instead. The limits check the formulas: as $t$ reaches the payment date, $P\to1$ and the next equity payment's value becomes the payment itself; just after a reset, $E_t=E_0$ and $A_t=1$, so each leg is worth $N(1-P_n)$ and the swap zero.

Hull's text gives the value of the next equity cash flow as $L(E-E_0)/E_0$ (p.~778), which omits discounting of the notional and is correct only at the payment date or with zero rates. His Technical Note~19, to which the passage refers, values the swap to the equity receiver at $LE/E_0-L(1+R_0\tau_0)/(1+R\tau)$, where $R$ is the rate from now to the next payment \citep[Technical Note~19]{hull-tn19}. Separating the floating payment $LR_0\tau_0/(1+R\tau)$ from it leaves $LE/E_0-L/(1+R\tau)$ for the equity payment, the corrected expression.

\qffield{Practice} \emph{Reproduced}: Hull's Examples~7.2 and~7.3, receiving 3 percent on 1{,}200 million yen and paying 4 percent on 10 million dollars for three years, with flat rates of 2.5 percent in dollars and 1.5 percent in yen and 110 yen per dollar: dollar bond 10.4191 million, yen bond 1{,}252.01 million, value 0.9628 million dollars by both methods. Hull prints 0.9629, but his own forward present values ($-0.0677$, $-0.0629$, $+1.0934$) and bond values both give 0.9628 (pp.~190--192).

\emph{Original} (curve of \hyperref[def:swaps-as-forward-strips]{swaps as forward strips}): a five-year receive-fixed swap at 3.75 percent on 100 million; a year later par rates are 4.10, 4.05, 4.15 and 4.25 percent for one to four years, so the swap is worth $-1.8079$ million by bonds and by forwards, $(3.75-4.25)$ percent times the four-year annuity times notional. A EUR/USD swap on EUR~100 million and USD~108 million at 1.08 with par coupons (2.50 percent EUR, 4.00 percent USD) is worth zero; if the euro rises to 1.12 dollars with rates unchanged, the party paying euros loses 4.0 million dollars. Entering 1.12 as euros per dollar instead would show an 18.7 million gain.

\emph{Equity swaps} (original, same curve). An equity-for-fixed swap has fair fixed rate 3.75 percent, the par swap rate. Half a year before the second reset, with the index up 6 percent since the last reset and a flat 4 percent curve: the next equity payment is worth $100\text{ million}\times(1.06-e^{-0.04\times0.5})=7.980133$ million, not the 6.000 million that $N(E_t/E_0-1)$ gives; the three later payments are worth 11.084044 million; the whole equity leg is 19.064 million, the fixed leg 13.861 million and the receive-equity side $+5.204$ million. Against compounded SOFR instead, with 1.900 percent realised since the reset ($A_t=1.019$), the floating leg is worth 3.880133 million for the next payment plus 11.084044 million for the later ones, and the receive-equity side $100\text{ million}\times(1.06-1.019)=4.100$ million.

\qffield{Conditions} One curve, a floater at par on reset dates, and no credit or funding adjustment: cash-collateralised trades are discounted at OIS, with term-rate projection in \hyperref[mod:multi-curve-fra-swap-valuation]{multi-curve valuation}. Hull's calculations ignore day counts and holiday calendars (p.~185). The equity swap formulas assume a total-return index (dividends reinvested), payment on the reset date, one curve, and no lookback or payment lag on the floating leg; a price-return index, a payment delay or a separate discounting curve each changes them.

\qffield{Pitfalls} Inverting a currency quote. Forgetting the principal exchanges in a currency swap, which carry most of its exchange-rate risk. Valuing the next equity payment as the return to date, $N(E_t/E_0-1)$, or confusing it with the whole equity leg or the net swap value. Treating swap value as the only exposure: Hull separates credit and market risk (sec.~7.11, p.~193).

\qffield{See also} \hyperref[def:swaps-as-forward-strips]{swaps as forward strips}; \hyperref[def:interest-rate-swaps]{interest-rate swaps}; \hyperref[mech:otc-clearing-collateral]{collateral}.

\end{qfmethod}
